% 原PDF161页附录4.F—165页附录4.G标题之前。
\section*{附录4.F\quad 动态结构方程模型参数的广义最小二乘估计量}
\addcontentsline{toc}{section}{附录4.F 动态结构方程模型参数的广义最小二乘估计量}
\markright{附录4.F 动态结构方程模型参数的广义最小二乘估计量}

定义矩阵
\[
\Gamma=\begin{bmatrix}
\Gamma^1&0&\cdots&0\\0&\Gamma^2&\cdots&0\\
\vdots&\vdots&\vdots&\vdots\\0&0&\cdots&\Gamma^T
\end{bmatrix},\quad
\Gamma_i=\begin{bmatrix}
\Gamma_i^1&0&\cdots&0\\0&\Gamma_i^2&\cdots&0\\
\vdots&\vdots&\vdots&\vdots\\0&0&\cdots&\Gamma_i^T
\end{bmatrix},
\]
\[
H=\begin{bmatrix}
H^1&0&\cdots&0\\0&H^2&\cdots&0\\
\vdots&\vdots&\vdots&\vdots\\0&0&\cdots&H^T
\end{bmatrix},\quad
H_i=\begin{bmatrix}
H_i^1&0&\cdots&0\\0&H_i^2&\cdots&0\\
\vdots&\vdots&\vdots&\vdots\\0&0&\cdots&H_i^T
\end{bmatrix}.
\]
其中 $i$ 表示第 $i$ 个方程；$\Gamma^t,\Gamma_i^t$ 定义于式（4.178），
$H^t=\begin{bmatrix}A^t\\B^t\end{bmatrix}$，
$A^t,B^t$ 分别定义于式（4.179）、（4.180）；$A_i^t,B_i^t,H_i^t$ 表示元素出现在第 $i$ 个方程中的矩阵。
式（4.196）定义了
\begin{equation}
\Pi^t=-H^t(\Gamma^t)^{-1},\tag{4.F.1}
\end{equation}
而矩阵 $\Pi$ 的定义见式（4.200）。在式（4.F.1）两侧右乘 $\Gamma^t$，得到
\begin{equation}
\Pi^t\Gamma^t=-H^t.\tag{4.F.2}
\end{equation}
注意
\begin{equation}
\begin{aligned}
\Pi\Gamma
&=\begin{bmatrix}\Pi^1&\cdots&0\\\vdots&\vdots&\vdots\\0&\cdots&\Pi^T\end{bmatrix}
\begin{bmatrix}
\Gamma^1&0&\cdots&0\\0&\Gamma^2&\cdots&0\\
\vdots&\vdots&\vdots&\vdots\\0&0&\cdots&\Gamma^T
\end{bmatrix}\\
&=\begin{bmatrix}
\Pi^1\Gamma^1&0&\cdots&0\\0&\Pi^2\Gamma^2&\cdots&0\\
\vdots&\vdots&\vdots&\vdots\\0&0&\cdots&\Pi^T\Gamma^T
\end{bmatrix}.
\end{aligned}\tag{4.F.3}
\end{equation}
将式（4.F.2）代入式（4.F.3），得到
\begin{equation}
\Pi\Gamma=-\begin{bmatrix}
H^1&0&\cdots&0\\0&H^2&\cdots&0\\
\vdots&\vdots&\vdots&\vdots\\0&0&\cdots&H^T
\end{bmatrix}=-H.\tag{4.F.4}
\end{equation}
类似地，
\begin{equation}
\begin{aligned}
\Pi\Gamma_i
&=\begin{bmatrix}
\Pi^1\Gamma_i^1&0&\cdots&0\\0&\Pi^2\Gamma_i^2&\cdots&0\\
\vdots&\vdots&\vdots&\vdots\\0&0&\cdots&\Pi^T\Gamma_i^T
\end{bmatrix}\\
&=-\begin{bmatrix}
H_i^1&0&\cdots&0\\0&H_i^2&\cdots&0\\
\vdots&\vdots&\vdots&\vdots\\0&0&\cdots&H_i^T
\end{bmatrix}=-H_i.
\end{aligned}\tag{4.F.5}
\end{equation}
将式（4.209）代入式（4.F.5），得到
\begin{equation}
\begin{aligned}
\Pi\Gamma_i
&=(W'W)^{-1}W'Y\Gamma_i\\
&=(W'W)^{-1}W'\begin{bmatrix}Y^1&\cdots&Y^T\end{bmatrix}
\begin{bmatrix}
\Gamma_i^1&0&\cdots&0\\0&\Gamma_i^2&\cdots&0\\
\vdots&\vdots&\vdots&\vdots\\0&0&\cdots&\Gamma_i^T
\end{bmatrix}\\
&=(W'W)^{-1}W'\begin{bmatrix}Y^1\Gamma_i^1&Y^2\Gamma_i^2&\cdots&Y^T\Gamma_i^T\end{bmatrix}
=-H_i.
\end{aligned}\tag{4.F.6}
\end{equation}
在式（4.F.7）两侧乘以 $W'W$，得到
\begin{equation}
W'\begin{bmatrix}Y^1\Gamma_i^1&Y^2\Gamma_i^2&\cdots&Y^T\Gamma_i^T\end{bmatrix}
=-W'WH_i.\tag{4.F.7}
\end{equation}
注意
\begin{equation}
Y^t\Gamma_i^t=Y^t\begin{bmatrix}-1\\\gamma_i^t\\0\end{bmatrix}
=-y_i^t+Y_i^t\gamma_i^t.\tag{4.F.8}
\end{equation}
接着计算 $WH_i$。由 $W,H_i$ 的定义，
\begin{equation}
\begin{aligned}
WH_i
&=\begin{bmatrix}W^1&\cdots&W^T\end{bmatrix}
\begin{bmatrix}
H_i^1&0&\cdots&0\\0&H_i^2&\cdots&0\\
\vdots&\vdots&\vdots&\vdots\\0&0&\cdots&H_i^T
\end{bmatrix}\\
&=\begin{bmatrix}W^1H_i^1&W^2H_i^2&\cdots&W^TH_i^T\end{bmatrix}.
\end{aligned}\tag{4.F.9}
\end{equation}
定义
\begin{equation}
W_i^t=\begin{bmatrix}Y_{li}^{t-1}&X_i\end{bmatrix}.\tag{4.F.10}
\end{equation}
于是
\begin{equation}
W^t=\begin{bmatrix}Y_l^{t-1}&X\end{bmatrix}
=\begin{bmatrix}Y_{li}^{t-1}&(Y_{li}^{t-1})^*&X_i&(X_i)^*\end{bmatrix}.\tag{4.F.11}
\end{equation}
回顾
\begin{equation}
H_i^t=\begin{bmatrix}\alpha_i^t\\0\\\beta_i^t\\0\end{bmatrix}.\tag{4.F.12}
\end{equation}
由式（4.188）、（4.194）、（4.195）、（4.F.10）、（4.F.11）及（4.F.12），
\begin{equation}
\begin{aligned}
W^tH_i^t
&=\begin{bmatrix}Y_{li}^{t-1}&(Y_{li}^{t-1})^*&X_i&(X_i)^*\end{bmatrix}
\begin{bmatrix}\alpha_i^t\\0\\\beta_i^t\\0\end{bmatrix}\\
&=Y_{li}^{t-1}\alpha_i^t+X_i\beta_i^t
=W_i^t\begin{bmatrix}\alpha_i^t\\\beta_i^t\end{bmatrix},
\quad t=1,\ldots,T.
\end{aligned}\tag{4.F.13}
\end{equation}
由式（4.F.7），
\begin{equation}
W'Y^t\Gamma_i^t=-W'WH_i^t.\tag{4.F.14}
\end{equation}
利用式（4.F.8），得到
\begin{equation}
W'Y^t\Gamma_i^t=-W'y_i^t+W'Y_i^t\gamma_i^t.\tag{4.F.15}
\end{equation}
结合式（4.F.13）、（4.F.14）、（4.F.15），得到
\[
-W'y_i^t+W'Y_i^t\gamma_i^t=-W'W_i^t\begin{bmatrix}\alpha_i^t\\\beta_i^t\end{bmatrix},
\]
从而
\begin{equation}
\begin{aligned}
W'y_i^t
&=W'Y_i^t\gamma_i^t+W'W_i^t\begin{bmatrix}\alpha_i^t\\\beta_i^t\end{bmatrix}\\
&=W'\left(Y_i^t\gamma_i^t+W_i^t\begin{bmatrix}\alpha_i^t\\\beta_i^t\end{bmatrix}\right)\\
&=W'\begin{bmatrix}Y_i^t&W_i^t\end{bmatrix}
 \begin{bmatrix}\gamma_i^t\\\alpha_i^t\\\beta_i^t\end{bmatrix}\\
&=W'Z_i^t\delta_i^t,\quad t=1,\ldots,T.
\end{aligned}\tag{4.F.16}
\end{equation}
回顾
\begin{equation}
y_i=\begin{bmatrix}y_i^1\\y_i^2\\\vdots\\y_i^T\end{bmatrix},\quad
Z_i\begin{bmatrix}
Z_i^1&0&\cdots&0\\0&Z_i^2&\cdots&0\\
\vdots&\vdots&\vdots&\vdots\\0&0&\cdots&Z_i^T
\end{bmatrix},\quad
\delta_i=\begin{bmatrix}\delta_i^1\\\delta_i^2\\\vdots\\\delta_i^T\end{bmatrix}.\tag{4.F.17}
\end{equation}
结合式（4.F.16）与式（4.F.17），得到
\begin{equation}
W_*'y_i=W_*'Z_i\delta_i.\tag{4.F.18}
\end{equation}
在式（4.189）两侧乘以 $W_*'$，得到
\begin{equation}
W_*'y_i=W_*'Z_i\delta_i+W_*'e_i.\tag{4.F.19}
\end{equation}
由式（4.190），$W_*'e_i$ 的协方差矩阵为
\begin{equation}
\Lambda_i=\operatorname{cov}(W_*'e_i,W_*'e_i)=W_*'V_iW_*,\tag{4.F.20}
\end{equation}
其中
\begin{equation}
V_i=\operatorname{cov}(e_i,e_i)=\Sigma_{ii}\otimes I_n,\tag{4.F.21}
\end{equation}
并且
\begin{equation}
\Sigma_{ii}=\begin{bmatrix}
\sigma_{ii}^{11}&\cdots&\sigma_{ii}^{1T}\\
\vdots&\vdots&\vdots\\\sigma_{ii}^{T1}&\cdots&\sigma_{ii}^{TT}
\end{bmatrix}.\tag{4.F.22}
\end{equation}
对模型（4.F.19）应用广义最小二乘，得到
\begin{equation}
\widehat\delta_i=
[(W_*'Z_i)'\Lambda_i^{-1}W_*'Z_i]^{-1}
(W_*'Z_i)'\Lambda_i^{-1}W_*'y_i.\tag{4.F.23}
\end{equation}
协方差矩阵 $\Sigma_{ii}$ 估计为
\begin{equation}
\widehat\Sigma_{ii}=\frac1{(n-m_1+1-l_i-k_i)T}
(y_i-Z_i\delta_i)(y_i-Z_i\delta_i)'.\tag{4.F.24}
\end{equation}
估计量 $\widehat\delta_i$ 的协方差矩阵为
\begin{equation}
\begin{aligned}
\Sigma_\delta=\operatorname{cov}(\widehat\delta_i,\widehat\delta_i)
={}&[(W_*'Z_i)'\Lambda_i^{-1}W_*'Z_i]^{-1}
(W_*'Z_i)'\Lambda_i^{-1}W_*'\operatorname{cov}(y_i,y_i)\\
&{}*W_*'\Lambda_i^{-1}(W_*'Z_i)
[(W_*'Z_i)'\Lambda_i^{-1}W_*'Z_i]^{-1}\\
={}&[(W_*'Z_i)'\Lambda_i^{-1}W_*'Z_i]^{-1}.
\end{aligned}\tag{4.F.25}
\end{equation}
